\centerline{\bf \Huge Лемма Холла и теорема Кёнига.}
\centerline{\bf \Huge }

\textbf{Определение.} \textit{Полным паросочетанием} в графе называется максимальное независимое множество ребер(то есть те ребра, у которых нет общих вершин).

Нас будут интересовать именно максимальные паросочетания, так как просто несколько независимых ребер найти не так уж и сложно. 

\textbf{Определение.} Набор ребер в графе называется \textit{рёберным покрытием}, если каждая вершина графа принадлежит хотя бы одному ребру этого набора.

\textbf{Определение.} Набор вершин в графе называется \textit{вершинным покрытием}, если каждое ребро графа принадлежит хотя бы одной из вершин этого набора.

Теперь сформулируем две очень важные теоремы, связывающие рёберные и вершинные покрытия с максимальным паросочетанием.

\problem{1} \textbf{Теорема Галлаи.} Дан граф на $n$ вершинах, среди которых нет изолированных. Сумма размеров максимального паросочетания и минимального реберного покрытия равна $n$.

\textbf{Доказательство:}
Рассмотрим максимальное паросочетание, пусть в нём $2k$ вершин. Любая из оставшихся вершин может быть соединена ребром лишь с вершинами паросочетания, причём из различных оставшихся вершин не могут вести рёбра в различные вершины пары (иначе ребро в паросочетании можно заменить на два). Так как в минимальном рёберном покрытии из каждой не входящей в паросочетание вершины выходит хотя бы по ребру, выделим сперва ровно по одному такому, их получилось $n - 2k$ рёбер. Теперь заметим, что для этого рёберного покрытия существует также минимальное рёберное покрытие, которое помимо выделенных $n- 2k$ рёбер содержит лишь рёбра паросочетания. Притом если оно не содержит ребро паросочетания, значит к обеим его вершинам ведут рёбра из вершин, не принадлежащих паросочетанию, а мы указывали, что такого быть не может. Таким образом, существует минимальное рёберное покрытие из $n - 2k$ рёбер и ещё $k$ рёбер паросочетания. Итого, сумма числа рёбер в максимальном паросочетании $k$ и числа рёбер в минимальном покрытии $n - k$ действительно равна $n$.

\newpage

\hspace{3mm}

\problem{2} \textbf{Теорема Кёнига.}

а) В прямоугольной таблице закрашено несколько клеток. Рядом будем называть столбец или строку. Докажите, что минимальное число рядов, которыми можно покрыть все закрашенные клетки, совпадает с максимальным числом не бьющих друг друга ладей, которые можно расставить в закрашенные клетки.

б) Дан двудольный граф. Вершинным покрытием будем называть любое такое множество
вершин $W$ , что у каждого ребра хотя бы один из концов принадлежит $W$. Паросочетание --- множество рёбер, не имеющих общих концов. Докажите, что число вершин минимального вершинного покрытия совпадает с числом рёбер в максимальном паросочетании.

Понятно, что здесь приведены две формулировки одной и той же задачи. Прежде чем доказывать теорему Кёнига, сформулируем и докажем лемму Холла, которая нам поможет.

\textbf{Лемма Холла.} Есть $n$ юношей и несколько девушек. Некоторые
юноши знают некоторых девушек. Известно, что для любых $k$ юношей (для всех $1 \leq k \leq n$) общее число знакомых им девушек не менее $k$. Тогда всех юношей можно поженить, каждого на знакомой ему девушке.

\textbf{Доказательство:} Доказывать лемму Холла будем с помощью так называемых "критических множеств"\,, но сами эти слова в ходе решения произносить не будем.

Будем рассуждать по индукции. При n = 1 теорема, очевидно, верна. Пусть она верна для любого числа юношей, меньшего n. Докажем, что она верна и для n юношей.

Могут быть только два случая:

\textbf{Случай 1.} При некотором $k \ (1 \leq k \leq n)$ найдётся группа из $k$ юношей, которые дружат (все вместе) ровно с $k$ девушками.

Обозначим множество этих юношей через Ю$_k$, а множество их подруг — через Д$_k$. Для группы юношей Ю$_k$ выполнено индукционное предположение. При этом их невесты исчерпывают всё множество Д$_k$ (в нём ровно $k$ девушек). Выбросим из рассмотрения этих юношей и девушек и докажем, что для оставшихся $n-k$ юношей выполнено условие теоремы, и, следовательно, можно воспользоваться предположением индукции, чтобы женить этих $n-k$ юношей.

Рассмотрим произвольную компанию из этих $n-k$, состоящую из $r$ человек. Докажем, что число их подруг (обозначим его $s$) больше либо равно, чем $r$. (Это не очевидно, потому как некоторые рёбра дружбы, возможно, вели к тем девушкам, которых мы удалили из рассмотрения.) Вернём назад и объединим этих $r$ юношей с группой Ю$_k$. В новой группе окажется $r + k$ юношей. Так как в множестве Д$_k$ содержится все подруги юношей из Ю$_k$ (и, возможно, ещё какие-то подруги юношей из второй компании), то у этих $k + r$ юношей в самом начале было $k + s$ подруг. По исходному условию $k + r \leq k + s$. Следовательно, $r \leq s.$ Значит, в этом случае теорема верна. 

\textbf{Случай 2.} Любая группа из $k$ юношей $(1 \leq k \leq n)$ дружит не менее чем с $k + 1$ девушкой.

В этом случае всё просто. Женим любого юношу на любой его подруге. Исключим эту пару из рассмотрения. Докажем, что остаётся $n-1$ юноша, для которых выполняется условие теоремы, и, следовательно, можно воспользоваться предположением индукции. Возьмём любую группу из $k - 1$ юноши. В самом начале у них было не менее $k+1$ подруг. Добавить очередного юношу - и он получил на одну ещё меньше, значит, группа имеет $k$ девушек. Это полностью доказывает теорему.

Теперь докажем теорему Кёнига.

\textbf{Доказательство теоремы Кёнига:}
Рассмотрим любое минимальное вершинное покрытие $A$. Порвём все рёбра между вершинами из покрытия (дальше работаем только с порванным графом). Докажем, что для вершин $L \subset A$ покрытия, лежащих в левой доле, выполнено условие леммы Холла. Предположим противное: пусть нашлось подмножество $X \subset L$, множество $Y$ друзей которого имеет меньшую мощность. Но тогда $(A \setminus X) \cup Y$ — вершинное покрытие с меньшим числом вершин (порванные рёбра покрыты с правой стороны, которую мы не трогали). Противоречие. Значит, условие леммы Холла выполнено и для $L$ можно найти паросочетание. Аналогично для вершин покрытия в правой доле. Построено паросочетание размера не меньше минимального вершинного покрытия, а значит максимальное паросочетание тем более не меньше.


\newpage

\hspace{3mm}

\textbf{\Large Примеры применения теоремы Кёнига.}

\textbf{Пример 1.}
\begin{enumerate}
    \item  Пусть $G$ --- двудольный граф, в каждой доле по $n$ вершин и в нём больше чем $(k - 1)n$   рёбер. Докажите, что в нём найдётся паросочетание, в котором хотя бы $k$ рёбер.
    \item Докажите, что из $51$   различного натурального числа, меньшего $100$,   можно выбрать $6$ чисел, попарно отличающихся в обоих разрядах (считаем, что у однозначных чисел в разряде десятков стоит $0$).
\end{enumerate}

\textbf{Доказательство:} 
\begin{enumerate}
    \item Предположим, что в любом паросочетании не более $k -1$   ребра. По теореме Кёнига имеем, что в этом графе есть вершинное покрытие $S$,   состоящее из не более $k - 1$   вершины, причём степень каждой вершины в $S$   не более $n$.   Следовательно, рёбер не более $n \cdot (k - 1)$,   что противоречит условию.
    \item Рассмотрим двудольный граф, вершины левой доли --- цифры в десятках, правой --- цифры в единицах. Рёбра --- выбранные числа. Далее применяем пункт (1).  
\end{enumerate}

\textbf{Пример 2.} В двудольном графе по $n$ вершин в каждой доле и степень каждой вершины не менее $\dfrac{n}{2}$. Докажите, что в нем есть паросочетание, содержащее все вершины.

\textbf{Доказательство:} Предположим, что такого паросочетания нет. Тогда по теореме Кёнига есть вершинное покрытие, содержащее не более $n - 1$ вершин. Следовательно, это вершинное покрытие содержит не более $\dfrac{n-1}{2}$ вершин в одной из долей. Но тогда это не вершинное покрытие, так как степень каждой вершины из другой доли больше $\dfrac{n}{2} > \dfrac{n-1}{2}$. Получили противоречие, значит есть паросочетание, содержащее все вершины.

\textbf{\Large Примеры применения леммы Холла.}

\textbf{Пример 1.} Дана таблица $n \times n$.   Её первые $k < n$   строчек заполнены натуральными числами от $1$ до $n$ так, что в каждой строке и в каждом столбце все числа различны. Докажите, что можно расставить натуральные числа от $1$ до $n$ в оставшиеся клетки таблицы так, чтобы по-прежнему выполнялось это условие.


\newpage

\hspace{3mm}

\textbf{Доказательство:} Достаточно показать, что мы можем заполнить ещё одну строчку с соблюдением условий. Рассмотрим двудольный граф, вершинами левой доли будут числа от $1$   до $n$,   вершинами правой доли --- столбцы, и проведём все рёбра от чисел к тем столбцам, в которых их нет. Заметим, что степени всех вершин равны по $n - k$.   Воспользуемся таким следствием из леммы Холла. Если в двудольном графе в одной доле $m$   вершин, а в другой не меньше $m$, причём степени вершин в каждой доле одинаковые между собой, то тогда можно выделить паросочетание в доле с $m$ вершинами. Следовательно, получаем, что в рассматриваемом графе есть паросочетание, покрывающее левую долю (а значит и правую). Тогда следуя этому паросочетанию мы можем расставить числа в $k + 1$   строке, а значит, и во всей таблице.

\textbf{Пример 2.} Имеются $27$   карточек с числами от $1$ до $27$.   Двое показывают следующий фокус. Первый получает $k$ карточек, выбранные случайным образом. Одну из них он убирает, а $k - 1$ оставшихся выкладывает в ряд. Второй должен назвать спрятанную карточку. Могут ли участники договориться так, чтобы по выложенным карточкам можно было определить спрятанную, если

(a) k= 4?  

(b) k= 14?  

\textbf{Доказательство:} Рассмотрим двудольный граф, в одной доле которого — сочетания из $27$ по $4$, в другой — размещения из $27$ по $3$, и ребро соединяет размещение с сочетанием, если все элементы размещения входят в сочетание.

Подсчитаем степень произвольной вершины из первой доли. Здесь нужно ответить на вопрос: сколько вариантов действия у первого участника фокуса? Он может убрать любую из четырёх карт, а три оставшиеся положить в ряд одним из $3!$ способов. Значит, степень вершины первой доли равна $4 \cdot 3! = 24$. Степень произвольной вершины второй доли равна количеству способов действия второго участника фокуса. Поскольку он может назвать любое из $27$ чисел, кроме тех трёх, которые он видит на карточках, и здесь степень вершины равна $24$. Тогда по лемме Холла существует паросочетание. Действительно, выберем из первой доли k вершин, тогда их суммарная степень равна $24k$. Пусть они соединены с $l$ вершинами из другой доли, тогда $24k \leq 24l \leq l \cdot l$. Аналогично во втором пункте. Давайте сделаем биекцию между наборами из 14 карт, которые видит первый, и наборами из $13$ карт который видит второй следующим образом. Рассмотрим двудольный граф, в одной доле которого — сочетания из $27$ по $14$, в другой — сочетания из $27$ по $13$, и ребро соединяет сочетания, если набор $13$ карт является подмножеством $14$. Тогда степени всех вершин в графе равны $14$, а ещё $C_{14}^{27} = C_{13}^{27}$, тогда по лемме Холла существует паросочетание.

\textbf{Пример 3.} Даны $n$ мальчиков и $2n - 1$   конфета. Докажите, что можно дать каждому мальчику по конфете так, чтобы мальчику, которому не нравится его конфета, не нравились и конфеты остальных мальчиков.

\textbf{Доказательство:} Рассмотрим двудольный граф, вершинами которого будут мальчики и конфеты, а рёбра будем проводить от мальчика к конфете, если та ему нравится. Если для доли мальчиков выполняется условие леммы Холла, то найдём паросочетание, покрывающее мальчиков и раздадим конфеты мальчикам, следуя ему. Предположим, что условие леммы Холла не выполняется. Пусть $A$ — это максимальное по включению множество мальчиков, для которого не выполняется лемма Холла, а $B$ — это множество конфет, которые нравятся мальчикам из A. Понятно, что $n = |A| > |B|$. Удалим из графа вершины из множеств $A$ и $B$. Так как $n \geq B$, то осталось не менее $n$ конфет. Для оставшегося графа для доли мальчиков выполнено условие леммы Холла, так как иначе множество $A$ было не максимальным по включению. Значит в оставшемся графе есть паросочетание, покрывающее мальчиков. Раздадим оставшимся мальчикам конфеты, следуя этому паросочетанию, а мальчикам из $A$ раздадим оставшиеся конфеты (их хватит, так как после удаления $B$ осталось не менее $n$ конфет). Нетрудно убедиться, что при такой раздаче конфет все условия задачи выполнены.
